\section{Preference Optimization：谁定义“更好的回答”}

Prompting 能改变一次推理，但若希望模型长期偏向有帮助、安全或符合特定风格的行为，就需要更新参数。Preference optimization 的共同输入是 chosen/rejected 或 reward signal；不同方法的差异在于是否训练 reward model、是否在线采样、如何构造 baseline，以及怎样表示人群偏好。

\subsection{RLHF：reward model 加 policy optimization}

RLHF 通常先拟合 reward model $r_{\phi}(x,y)$，再优化
\[
\max_{\theta}\;\mathbb{E}_{y\sim\pi_{\theta}(\cdot\mid x)}[r_{\phi}(x,y)]
-\beta D_{\mathrm{KL}}(\pi_{\theta}\|\pi_{\mathrm{ref}}).
\]
KL 项限制 policy 偏离 reference model。系统成本来自 preference collection、reward model、rollout 与 RL 稳定性；reward hacking 则说明代理分数不等于真实用户价值。


\lecturefigure{slide-078.jpg}{RLHF pipeline：人类比较、reward model 与 policy optimization}{CS25 V5 official deck, p.~78}

\teachervoice{00:29:13--00:31:32，讲者把 prompting 与 preference training 分开：前者在 inference time 提供 scaffold，后者用 human/AI feedback 改变模型参数。这个边界比 acronyms 更重要。}

\subsection{DPO 与 RLAIF：改变训练信号来源}

DPO 直接用 preference pairs 优化 policy 相对 reference 的 log-ratio：
\[
\mathcal{L}_{\mathrm{DPO}}=-\log\sigma\!\left(\beta\left[
\log\frac{\pi_{\theta}(y^+\mid x)}{\pi_{\mathrm{ref}}(y^+\mid x)}-
\log\frac{\pi_{\theta}(y^-\mid x)}{\pi_{\mathrm{ref}}(y^-\mid x)}
\right]\right).
\]
它省去显式 reward-model RL loop，但仍依赖 pair quality、reference 与偏好分布。RLAIF 用 AI judge 替代部分人类标注，降低成本，却把 judge bias、能力和可操纵性引入监督链。


\lecturefigure{slide-079.jpg}{DPO：从 preference pair 直接优化 policy}{CS25 V5 official deck, p.~79}


\lecturefigure{slide-080.jpg}{RLAIF：用 AI feedback 扩展偏好标注}{CS25 V5 official deck, p.~80}

\begin{warningbox}{Judge capability 是训练上限之一}
如果 AI judge 无法识别事实错误、长程策略或细微安全风险，它会把噪声写入 preference data。用更强 judge 不自动解决同质偏差；需要 human audit、多 judge disagreement、adversarial prompts 和 outcome-based checks。
\end{warningbox}

\subsection{GRPO、KTO 与 heterogeneous preferences}

GRPO 用一组 sampled responses 的相对回报构造 advantage，避免独立 value model 的部分成本。简化表示为
\[
A_i=\frac{r_i-\mu_G}{\sigma_G+\epsilon},
\]
其中 $\mu_G,\sigma_G$ 是同组 reward 的均值和标准差。它依赖 group diversity 与 reward reliability。


\lecturefigure{slide-081.jpg}{GRPO：组内相对 reward 与 policy 更新}{CS25 V5 official deck, p.~81}

KTO 从 prospect theory 借用 loss aversion，允许只有 desirable/undesirable 标签而非严格 pairs；它强调负面结果的非对称价值。Variational Preference Learning 则把不同 demographic 或 profile 的偏好建模为潜变量，承认“一个平均 reward”可能抹平真实分歧。


\lecturefigure{slide-082.jpg}{KTO：loss aversion 与非配对偏好信号}{CS25 V5 official deck, p.~82}


\lecturefigure{slide-083.jpg}{Variational preference learning：个性化与异质人群偏好}{CS25 V5 official deck, p.~83}

\begin{warningbox}{Preference method 上线前的五层验证}
第一层检查 preference labels 的一致性、群体覆盖与敏感属性；第二层做 held-out pair accuracy，但不能把 reward-model accuracy 当最终产品质量；第三层用多维 benchmark 分开测 helpfulness、truthfulness、harmlessness、style 与 diversity；第四层在 live-like prompts 上做 blinded human comparison，并记录 disagreement 而非强行平均；第五层开展 adversarial evaluation，寻找 reward hacking、sycophancy、over-refusal 和 judge manipulation。还要比较 base/reference drift、KL、回答长度和校准。一个方法离线 win-rate 更高，若代价是模式坍缩、特定群体偏好被抹平或事实性下降，就不是总体进步。Slides 的 loss diagrams 只定义优化路径，验收必须回到真实行为分布。
\end{warningbox}

\begin{knowledgebox}{术语消化：reward、preference 与 value}
\textbf{Preference label} 表示两个回答的相对选择；\textbf{reward model} 把输入回答映射为标量；\textbf{value/critic} 估计未来回报；\textbf{reference policy} 约束更新不偏离原模型；\textbf{judge} 可以是人或 AI。RLHF、DPO、GRPO、KTO 的区别来自这些对象如何组合。读 slides 78--83 时先画出监督图，再看 loss；只背公式容易忽略 reward distribution、标注群体和 KL constraint 才是行为边界。
\end{knowledgebox}

\teachervoice{00:31:32--00:34:13，课堂把 GRPO、KTO 与 variational preferences 分别连接到 group-relative reward、loss aversion 和 demographic heterogeneity。老师强调的是 objective assumptions，而不是宣布某方法普遍最优。}

\begin{table}[H]
\centering
\small
\begin{tabularx}{\textwidth}{L{0.13\textwidth} L{0.20\textwidth} X X}
\toprule
方法 & 监督来源 & 省掉/新增的组件 & 关键风险 \\
\midrule
RLHF & 人类偏好 & reward model + RL & reward hacking、昂贵 rollout \\
DPO & pair data & 无显式 RL loop & pair bias、reference 依赖 \\
RLAIF & AI judge & 扩展标注吞吐 & judge bias 与同质错误 \\
GRPO & group rewards & 相对 baseline & group/reward 质量 \\
KTO & binary desirability & 不要求 pairs & value model 假设 \\
VPL & profile-conditioned data & 偏好潜变量 & 隐私、分群和公平性 \\
\bottomrule
\end{tabularx}
\caption{Preference optimization 方法的监督接口与失败模式。}
\end{table}

\begin{lstlisting}[caption={Preference data pipeline 的最小审计骨架}]
prompts = stratified_sample(real_use, safety_edges, long_tail)
responses = generate_candidates(policy, prompts)
labels = collect_preferences(humans, ai_judges, outcomes)
audit_disagreement(labels, by_domain=True, by_group=True)
train_preference_method(labels, reference_model)
validate(helpfulness, factuality, safety, calibration, diversity)
\end{lstlisting}

\subsection{本章小结}

Preference optimization 的核心不是 acronym，而是谁给信号、信号怎样转成 loss、模型相对什么 baseline 更新。RLHF、DPO、RLAIF、GRPO、KTO 和 VPL 在成本与假设上不同；任何方法都不能用单一 reward 替代真实多维产品目标。

